\section{The Arguments, Each at Its Strongest}

An argument about discipline is worth having only if both sides are stated in the form their best advocates would recognize. What follows states them that way and adjudicates neither, because the Church has not, and because this article's competence is the law rather than the prudence of the law.

\subsection{What is not in dispute}

Four things are settled by the texts already read, and no serious position on either side denies them.

Chastity binds everyone, in the form proper to each state of life, and is not what any of this is about. Marriage after ordination is excluded in both codes---CIC canon 1087, CCEO canon 804---and no party to the modern argument proposes changing that. The Eastern tradition of a married presbyterate is legitimate, ancient, and honoured by name in the Latin Church's own conciliar teaching (\emph{Presbyterorum ordinis} 16) and in the Eastern Code (CCEO c.~373). And the Latin obligation is ecclesiastical law: not divine law, not an entailment of the sacrament, and therefore---as a matter of possibility, whatever one thinks of the prudence---changeable by the legislator who made it.

A position that has to deny any of those four is arguing against the record.

\subsection{The case for the Latin discipline}

\textbf{Configuration to Christ.} The strongest form of this argument is not that a celibate priest is holier but that the priesthood is a sacramental representation, and that Christ's own state was celibate. \emph{Sacerdotalis caelibatus} 21: ``Christ remained throughout His whole life in the state of celibacy, which signified His total dedication to the service of God and men.'' The claim is about signification, not about merit, and it is why the argument is not answered by pointing out that married men can be holy---which no one disputes.

\textbf{The spousal sign.} \emph{Presbyterorum ordinis} 16 draws the image out: priests ``profess themselves before men as willing to be dedicated to the office committed to them---namely, to commit themselves faithfully to one man and to show themselves as a chaste virgin for Christ.'' \emph{Pastores dabo vobis} 29 makes it explicit that virginity ``makes evident, even in the renunciation of marriage, the `nuptial meaning' of the body.'' On this reading celibacy is not the absence of a relation but the visibility of one.

\textbf{The eschatological sign.} \emph{Presbyterorum ordinis} 16 again: priests ``give a living sign of the world to come, by a faith and charity already made present, in which the children of the resurrection neither marry nor take wives.'' The force of this argument is that a sign of the age to come has to be legible now, and that a discipline making celibacy optional would leave it invisible as a public sign of the priesthood, whatever it remained for individuals.

\textbf{Pastoral availability, in the canon's own words.} Canon 277 §1 gives its own reason---\latin{indiviso corde}, an undivided heart, and \latin{liberius}, more freely---and the reason is practical as well as theological: a man without a family has obligations that a man with one does not. \emph{Presbyterorum ordinis} 16's fuller version is that celibacy makes priests apt ``to accept, in a broad sense, paternity in Christ.''

\textbf{The antiquity argument.} A body of twentieth-century scholarship argues that the Western legislation of the fourth century did not invent an obligation but codified an apostolic-age discipline of perfect continence binding married clergy, so that the Latin discipline is continuous with the earliest Church and the divergence is Eastern rather than Western. This is the strongest historical form of the case, because it shifts the burden: if it is right, the Latin rule is a preservation rather than a medieval imposition. It is a live and contested thesis, and section~10 sets out both the evidence it rests on and the reasons its opponents give.\endnote{Reported at literature level. The principal modern statements of the apostolic-continence thesis are C.~Cochini, \emph{Origines apostoliques du célibat sacerdotal} (Paris: Lethielleux, 1981; English, San Francisco: Ignatius, 1990); A.~M. Stickler, \emph{The Case for Clerical Celibacy} (San Francisco: Ignatius, 1995); R.~Cholij, \emph{Clerical Celibacy in East and West} (Leominster: Fowler Wright, 1988); and S.~Heid, \emph{Zölibat in der frühen Kirche} (Paderborn: Schöningh, 1997; English, San Francisco: Ignatius, 2000). The principal statement of the contrary position is R.~Gryson, \emph{Les origines du célibat ecclésiastique du premier au septième siècle} (Gembloux: Duculot, 1970). None of these works was consulted at source for this article; they are named so that a reader can find the debate, and no claim here rests on any of them.}

\textbf{The Church's freedom.} Finally, and often left unsaid: the Church has the right to determine the conditions under which she confers orders, and the fact that a norm is not required by the nature of the sacrament is not an argument that it may not be imposed. \emph{Sacerdotalis caelibatus} 14 puts it as a positive judgment---``the present law of celibacy should today continue to be linked to the ecclesiastical ministry''---and a legislator's prudential judgment is not answered by showing that he was not compelled to make it.

\subsection{The case for change}

\textbf{It is law, and law serves.} The strongest form of this argument concedes everything in the previous subsection and observes that a law of fittingness is a means. If a means at some time and place obstructs the end---in the Amazon synod's formulation at n.~110, ``a right of the community to the celebration'' that ``derives from the essence of the Eucharist''---then the means is reconsidered. That is not an argument against celibacy but an argument about what a disciplinary norm is for.

\textbf{The Eastern witness is not an exception.} If a married presbyterate were a concession to weakness, its position in the Eastern Code would be different from what section~7 found: the state of married clerics is \latin{in honore habendus}, warranted by the practice of the primitive Church, in a canon whose other clause honours celibacy. A discipline honoured in one part of the Catholic Church cannot be called unsuitable in the other without saying something about the East that no Catholic document says.

\textbf{The principle is already conceded.} The Latin Church ordains married men to the diaconate as a matter of general law (c.~1042 1°) and to the presbyterate by reserved dispensation and papal derogation (c.~1047 §2 3°; \emph{Anglicanorum coetibus} VI §2). She recognizes married Eastern presbyters, and since 2014 has restored to Eastern Hierarchs outside the traditional territories the faculty to ordain them. What is at issue, on this view, is therefore not whether a married priesthood is compatible with Latin ecclesial life but how widely the existing provisions run.

\textbf{Charism and law.} \emph{Sacerdotalis caelibatus}, canon 277 §1, and \emph{Optatam totius} 10 all call celibacy a gift---\latin{peculiare Dei donum}, \latin{pretiosum donum Dei}. A gift is given to whom it is given. The argument is that a law requiring the gift as a condition of orders risks two failures at once: excluding men who have the vocation to priesthood but not to celibacy, and admitting men who accept celibacy as a price. \emph{Pastores dabo vobis} 50 makes the second half of that worry the Church's own: celibacy ``should not be considered just as a legal norm or as a totally external condition for admission to ordination.''

\textbf{The counter-historiography.} Against the antiquity argument stands the reading on which the New Testament and the earliest centuries attest married clergy without any general obligation of continence, the fourth-century Western canons are innovations rather than codifications, and the Eastern settlement at Trullo preserved the older practice. Section~10 sets out what each of the disputed texts actually says.

\subsection{Three arguments that do not work}

\emph{``Priests could marry until 1139.''} This misdescribes the law in a way that damages the side that uses it. What the twelfth-century legislation did was change the effect of an attempted marriage, not remove a permission; the prohibition is far older, and section~10 traces it. An argument for change resting on a misstatement of when the rule began invites a rebuttal on that point alone.

\emph{``The Eastern Churches prove the Latin rule is wrong.''} They prove that it is not necessary, which \emph{Presbyterorum ordinis} 16 had already conceded and which no defender of the discipline needs to deny. Between the necessary and the wrong lies the whole field of prudential legislation, which is where this dispute actually is.

\emph{``Celibacy causes abuse''---and its mirror image, that raising the question is an attack on the priesthood.} The law's own judgment about the classification of these delicts was stated in section~5: in Book~VI as revised in 2021, the abuse of a minor stands under Title~VI, offences against human life, dignity, and freedom, beside homicide and abduction, and not under Title~V, offences against special obligations, where clerical concubinage and attempted marriage stand. That is a legislative classification of what the delict is against, and it is the only thing here within this article's competence to report. Whether celibacy bears any causal relation to the incidence of abuse is an empirical question about which this article examined no evidence, states no conclusion, and offers no reassurance in either direction; it is answered, if at all, by research and not by canon law. What can be said is that the classification cuts against treating the crime as a variety of unchastity, and that a discipline argument which uses victims as ammunition on either side has stopped being a discipline argument.

\subsection{What would actually have to change, and what would not}

This subsection is the article's own synthesis, offered as a way of seeing the size of the question rather than as a proposal.

If the Latin Church decided to ordain married men to the presbyterate as a general matter, three norms would carry the change. Canon 277 §1 would have to be amended, or derogated from generally rather than case by case, since it is the source of the obligation. Canon 1042 1° would have to be widened, since it is the impediment. Canon 1037 would have to be adjusted, since it presently requires the public assumption of celibacy from every presbyteral candidate before the diaconate. Nothing else in the apparatus of sections~4 through~6 is strictly necessary to the change.

And three things would not change, on any proposal actually made by anyone. Canon 1087 would stand: a man ordained unmarried could still not marry afterwards, and a widowed priest could still not remarry. The episcopate would remain unmarried, as it is in every Eastern Church by canon 180 3° and in the Latin Church by the general obligation. And the two obligations of canon 277 §1 would have to be separated explicitly rather than silently, since the question the Latin Code has already left unanswered for the married permanent deacon (section~9) would arise for every married presbyter on the day the change took effect.

That last point is why this article has insisted on the distinctions from its first page. A Church that has not settled whether \latin{perfecta perpetuaque continentia} binds a married deacon has not settled the prior question that any change to the celibacy law would put at the centre of its own discipline.
